\chapter*{Summary of the thesis}
\addcontentsline{toc}{chapter}{Summary of the thesis}

\section*{Preface}
The evidence for dark matter is overwhelming, and yet it tells us almost nothing about what dark matter is made of. Galactic rotation curves stay flat where the visible mass predicts a decline; galaxies inside clusters move too fast to remain bound by the luminous matter alone; light from distant sources is bent by mass we cannot see; and the pattern of anisotropies in the cosmic microwave background is reproduced only when a non-luminous component is added to the budget. Each of these measurements, drawn from a different scale and a different epoch, points to the same conclusion: a form of matter that outweighs ordinary baryons by roughly a factor of five and interacts with light only through gravity. Within the standard cosmological model this component accounts for about 85\% of the matter density, cold and non-baryonic. The agreement across independent probes is remarkable, but it is also frustrating, because every one of these probes is gravitational. They tell us how much dark matter there is and where it sits; they say nothing about its particle nature. Proposed candidates span nearly ninety orders of magnitude in mass, from ultralight axions to primordial black holes, and the data alone single out none of them.

Among these candidates, weakly interacting massive particles (WIMPs) have long held a privileged place. A stable particle with an electroweak-scale mass and coupling, produced thermally in the hot early Universe, freezes out with a relic abundance that lands naturally near the observed dark matter density, a coincidence known as the WIMP miracle. The annihilation process that fixes this abundance does not switch off once the Universe cools: wherever dark matter piles up densely today, WIMPs should continue to annihilate, converting their rest mass into Standard Model particles such as gamma rays, neutrinos, positrons, and antiprotons. If dark matter is instead unstable, it may decay into the same states. Of the messengers this produces, gamma rays are the most informative. They travel undeflected by magnetic fields, so they point back to where they were made, and they carry the spectral fingerprint of the interaction that produced them. A gamma-ray signal from a region of concentrated dark matter would therefore be one of the cleanest handles we have on the particle behind the gravitational anomalies.

For more than fifteen years the Large Area Telescope (LAT) on board the Fermi Gamma-ray Space Telescope has surveyed the GeV sky where such signatures are expected, and its measurements sit at the core of this thesis. That long exposure has changed the character of the problem. The brightest targets have already been studied at length, and the signals that remain are faint, blended with astrophysical foregrounds, or hidden below the detection threshold. Standard methods classify the brightest objects and map the prominent diffuse components reliably, but they struggle with overlapping populations and with instrumental limits such as finite angular resolution and the detection threshold itself. Setting a threshold and counting what rises above it is a natural first move, yet it discards most of the information: a large fraction of the emitting population is left unresolved, unclassified, or simply unseen.

We argue in this thesis that progress in indirect dark matter searches now depends on statistical and machine-learning methods able to recover signals in these noise-dominated regimes, where simple thresholding fails. The work follows an arc from the Galactic Center outward to the cosmic web, tracking the expected signal into regimes of steadily lower signal-to-noise. As the signal weakens, the methodology has to grow more sophisticated to keep pace, and we draw in turn on traditional statistical inference, generative modeling, implicit-likelihood techniques, and deep learning, applied both to Fermi-LAT data and to forecasts for future instruments.

Five studies make up this progression. We begin with a measurement of the gamma-ray luminosity function of millisecond pulsars (MSPs) in Milky Way globular clusters, used to test whether an unresolved pulsar population can account for the Galactic Center Excess. We then search for dark matter subhalos among the unassociated Fermi-LAT sources, combining a generative mixture model with quantification learning. Moving below the catalog threshold, we reconstruct the source-count distribution of faint extragalactic sources with a deep neural network, and we build a probabilistic cataloging framework that extends gamma-ray catalogs into the sub-threshold regime, released publicly as the gamma-ray Photon Count Statistics (gPCS) catalog, a Python package accompanied by a FITS table. Finally, we forecast the sensitivity of the Cherenkov Telescope Array Observatory to the cross-correlation between gamma rays and galaxy catalogs, carrying the search out to the largest scales and the weakest expected signals.

\section*{The dark matter problem}
% <<DM-PROBLEM>>

\section*{The gamma-ray sky and the \textit{Fermi} Large Area Telescope}
% <<GAMMA-SKY>>

\section*{Objectives}
% <<OBJECTIVES>>

\section*{Methodology: statistics and machine learning for faint signals}
% <<METHODOLOGY>>

\section*{Results}
% <<RESULTS-MSP>>
% <<RESULTS-SUBHALOS>>
% <<RESULTS-DNDS>>
% <<RESULTS-CATALOG>>
% <<RESULTS-XCORR>>

\section*{Conclusions and future directions}
% <<CONCLUSIONS>>
